\documentclass[12pt,oneside]{book}

% ==========================================================
% METADATOS
% ==========================================================
\newcommand{\universidad}{Universidad de Buenos Aires (UBA)}
\newcommand{\facultad}{Facultad de Derecho}
\newcommand{\materia}{Derechos Reales}
\newcommand{\titulo}{Teoría General: Adquisición, Transmisión y Extinción de los Derechos Reales}
\newcommand{\autor}{Isaac Edgar Camacho Ocampo}
\newcommand{\anio}{2020}

% ==========================================================
% PAQUETES
% ==========================================================
\usepackage[utf8]{inputenc}
\usepackage[T1]{fontenc}
\usepackage[spanish,es-nodecimaldot]{babel}

\usepackage[
    top=2.5cm,
    bottom=2.5cm,
    left=3cm,
    right=2.5cm,
    headheight=15pt
]{geometry}

\usepackage{amsmath}
\usepackage{amssymb}
\usepackage{mathtools}

\usepackage{graphicx}
\usepackage{float}
\usepackage{caption}
\usepackage{subcaption}

\usepackage{booktabs}
\usepackage{array}
\usepackage{longtable}
\usepackage{tabularx}

\usepackage{enumitem}

\usepackage{xcolor}
\usepackage[most]{tcolorbox}

\usepackage{fancyhdr}
\usepackage{ragged2e}

\usepackage[
    colorlinks=true,
    linkcolor=blue!50!black,
    citecolor=green!40!black,
    urlcolor=blue!60!black,
    pdftitle={Teoría General: Adquisición, Transmisión y Extinción de los Derechos Reales},
    pdfauthor={Isaac Edgar Camacho Ocampo},
    pdfsubject={Apuntes universitarios},
    bookmarks=true
]{hyperref}

% ==========================================================
% CONFIGURACIÓN
% ==========================================================
\graphicspath{
    {./img/}
    {./graficos/}
    {./figuras/}
}

\setcounter{tocdepth}{2}
\setcounter{secnumdepth}{3}

\setlength{\parindent}{0pt}
\setlength{\parskip}{0.5em}

\setlist[itemize]{
    topsep=0.3em,
    itemsep=0.2em
}

\setlist[enumerate]{
    topsep=0.3em,
    itemsep=0.2em
}

\clubpenalty=10000
\widowpenalty=10000

\pagestyle{fancy}
\fancyhf{}
\fancyhead[L]{\nouppercase{\leftmark}}
\fancyhead[R]{\materia}
\fancyfoot[C]{\thepage}
\renewcommand{\headrulewidth}{0.4pt}
\renewcommand{\footrulewidth}{0pt}

\fancypagestyle{plain}{
    \fancyhf{}
    \fancyfoot[C]{\thepage}
    \renewcommand{\headrulewidth}{0pt}
}

% ==========================================================
% COMANDOS
% ==========================================================
\newcommand{\abs}[1]{\left|#1\right|}
\newcommand{\norm}[1]{\left\lVert#1\right\rVert}
\newcommand{\vect}[1]{\mathbf{#1}}
\newcommand{\deriv}[2]{\frac{d#1}{d#2}}
\newcommand{\pderiv}[2]{\frac{\partial#1}{\partial#2}}

% ==========================================================
% ENTORNOS / CAJAS
% ==========================================================
\newtcolorbox{definition}[1]{
    enhanced,
    breakable,
    title={Definición: #1},
    fonttitle=\bfseries,
    colback=blue!3,
    colframe=blue!50!black,
    boxrule=0.8pt,
    arc=2mm
}

\newtcolorbox{important}{
    enhanced,
    breakable,
    title={Importante},
    fonttitle=\bfseries,
    colback=yellow!8,
    colframe=orange!70!black,
    boxrule=0.8pt,
    arc=2mm
}

\newtcolorbox{example}[1]{
    enhanced,
    breakable,
    title={Ejemplo: #1},
    fonttitle=\bfseries,
    colback=green!4,
    colframe=green!40!black,
    boxrule=0.8pt,
    arc=2mm
}

\newtcolorbox{formula}[1]{
    enhanced,
    breakable,
    title={Fórmula / Regla: #1},
    fonttitle=\bfseries,
    colback=purple!4,
    colframe=purple!50!black,
    boxrule=0.8pt,
    arc=2mm
}

\newtcolorbox{exercise}[1]{
    enhanced,
    breakable,
    title={Ejercicio: #1},
    fonttitle=\bfseries,
    colback=gray!5,
    colframe=gray!60!black,
    boxrule=0.8pt,
    arc=2mm
}

\newtcolorbox{note}{
    enhanced,
    breakable,
    title={Nota},
    fonttitle=\bfseries,
    colback=white,
    colframe=gray!50,
    boxrule=0.6pt,
    arc=2mm
}

% ==========================================================
% DOCUMENTO
% ==========================================================
\begin{document}

% ----------------------------------------------------------
% PORTADA
% ----------------------------------------------------------
\begin{titlepage}

\thispagestyle{empty}

\begin{center}

\vspace*{1cm}

{\large \textsc{\universidad}}

\vspace{0.2cm}

{\large \textsc{\facultad}}

\vspace{3cm}

{\Huge \textbf{\materia}}

\vspace{1cm}

{\LARGE \titulo}

\vspace{0.4cm}

{\large Clase N.º 2}

\vspace{0.8cm}

\textit{\small Estudiar con constancia es construir, concepto a concepto, un conocimiento propio.}

\vspace{1.2cm}

\rule{0.70\textwidth}{0.5pt}

\vspace{0.5cm}

{\large Apuntes de estudio}

\vspace{0.5cm}

\rule{0.70\textwidth}{0.5pt}

\vfill

{\large \autor}

\vspace{0.3cm}

{\large \anio}

\end{center}

\vfill

\begin{flushleft}
\textit{Este es un modesto aporte a los estudiantes de ``\materia'' no pretende sustituir el material oficial de la materia.}
\end{flushleft}

\end{titlepage}

% ----------------------------------------------------------
% ÍNDICE
% ----------------------------------------------------------
\tableofcontents
\clearpage

% ==========================================================
% CAPÍTULO 1
% ==========================================================
\chapter[Caracteres y fundamento de los derechos reales]{Caracteres y fundamento de los derechos reales}

Los apuntes parten de un repaso de lo ya estudiado, que permite distinguir con nitidez las dos grandes esferas del derecho patrimonial: los \textbf{derechos reales} y los \textbf{derechos personales}, así como los caracteres propios de cada una. El tema de esta clase es la \textbf{teoría general} de los derechos reales, que se complementa teniendo en cuenta lo ya avanzado.

\section{Derechos reales y derechos personales}

\subsection{Caracteres propios de cada esfera}

Los derechos reales y los derechos personales constituyen dos esferas de derechos con caracteres propios, claramente diferenciables entre sí.

\subsection{El número cerrado de derechos reales (numerus clausus)}

Un carácter esencial de los derechos reales es su \textbf{número limitado y taxativo}: los derechos reales son \textbf{14 y solo 14}. A diferencia de los derechos personales, los derechos reales son \textbf{creados exclusivamente por la ley}, y hoy se encuentran todos codificados.

La competencia para crearlos es exclusivamente federal: se trata de \textbf{ley federal}, porque en materia de derechos reales solo la Nación es competente para legislar.

\subsection{Las normas de orden público en materia de derechos reales}

Los particulares deben ceñirse estrictamente a la regulación legal de cada derecho real, porque las normas que lo regulan son de \textbf{orden público}, expresan la voluntad de la Nación y, en consecuencia, no pueden ser modificadas por los particulares.

La razón de fondo es que el derecho real responde a un esquema de sociedad único para toda la Nación. Tiene que ver con el modo en que una Nación organiza su sistema de propiedad y su sistema de defensa de las relaciones de poder respecto de las cosas; con el modo en que una sociedad define el alcance de las facultades que está dispuesta a reconocer a cada titular de un derecho real; y con las limitaciones que está dispuesta a aceptar en función de los intereses de la comunidad.

\begin{formula}{Art. 12 del Código Civil y Comercial (CCyCN)}
Las convenciones particulares no pueden dejar sin efecto las leyes en cuya observancia está interesado el orden público.
\end{formula}

\section{Fundamento constitucional e internacional del derecho real}

\subsection{La constitucionalización del derecho privado}

El Código Civil y Comercial introduce un cambio de paradigma: se incorporan los tratados internacionales. Por ello puede hablarse de la \textbf{constitucionalización} y de la \textbf{internacionalización del derecho privado}, ya que los tratados de derechos humanos son hoy ley.

Los arts. 1 y 2 del CCyCN son fundamentos del Código y, a la vez, una \textbf{regla de interpretación}: todo el sistema es mirado bajo otra perspectiva.

\subsection{El sistema de propiedad privada y sus límites}

Conviven en tensión varias normas del sistema, según se resume en el siguiente cuadro.

% TODO: contenido incompleto o ambiguo en las notas originales (la numeración de los arts. 14 y 15 CCyCN y sus contenidos se conservan tal como figuran en las notas)
\begin{center}
\small
\renewcommand{\arraystretch}{1.3}
\begin{tabularx}{\textwidth}{@{}>{\raggedright\arraybackslash}p{0.17\textwidth}>{\raggedright\arraybackslash}X>{\raggedright\arraybackslash}p{0.27\textwidth}@{}}
\toprule
\textbf{Artículo} & \textbf{Contenido} & \textbf{Efecto} \\
\midrule
Art. 15 CCyCN & Las personas son las titulares de los derechos individuales sobre todos los bienes que integran su patrimonio. & Consagra el sistema de propiedad privada. \\
Art. 14 CCyCN & La ley no ampara el ejercicio abusivo de los derechos individuales cuando puede afectar al ambiente y a los derechos de incidencia colectiva. & Limita los derechos individuales frente al interés colectivo. \\
Art. 240 CCyCN & Establece un límite a los derechos individuales cuando resulte necesario compatibilizarlos con los derechos de incidencia colectiva. & Preeminencia del interés colectivo sobre el individual. \\
\bottomrule
\end{tabularx}
\end{center}

En el sistema jurídico existe una clara \textbf{preeminencia del interés colectivo por sobre el interés individual}. Por eso, en ocasiones (como en una situación de contingencia), resulta necesario modelar el alcance del ejercicio de un derecho real en función de los intereses de la comunidad.

\subsection{Derechos individuales y derechos de incidencia colectiva}

El derecho al ambiente y el derecho a respirar son ejemplos de derechos de incidencia colectiva que pueden entrar en tensión con el ejercicio de los derechos reales.

El derecho real se vincula con los grandes temas del derecho: el ejercicio del derecho de propiedad, los límites al dominio, el reconocimiento de la propiedad como un derecho humano y la función social. Se trata de temas susceptibles de enfrentar ideológicamente a distintos sectores de la doctrina nacional e internacional.

\section{Cuadro Cornell del capítulo}

\begin{longtable}{@{}>{\raggedright\arraybackslash}p{0.27\textwidth}>{\raggedright\arraybackslash}p{0.67\textwidth}@{}}
\toprule
\textbf{Preguntas / palabras clave} & \textbf{Notas / respuestas} \\
\midrule
\endfirsthead
\toprule
\textbf{Preguntas / palabras clave} & \textbf{Notas / respuestas} \\
\midrule
\endhead
\bottomrule
\endfoot

\textbf{¿Cuántos derechos reales existen y por qué?} &
Son 14 y solo 14. A diferencia de los derechos personales, los derechos reales son creados exclusivamente por la ley (numerus clausus), y todos están codificados. Solo la Nación es competente para legislarlos (ley federal). \\[0.8em]

\textbf{¿Por qué los particulares no pueden modificar las normas de los derechos reales?} &
Porque son normas de orden público que expresan la voluntad de la Nación y responden a un esquema de sociedad único para toda ella: el modo en que se organiza el sistema de propiedad, el alcance de las facultades reconocidas a cada titular y las limitaciones aceptadas en función de los intereses de la comunidad (art. 12 CCyCN). \\[0.8em]

\textbf{¿Qué es la constitucionalización del derecho privado?} &
Es la incorporación al Código Civil y Comercial de los tratados internacionales de derechos humanos, que hoy son ley. Los arts. 1 y 2 CCyCN son fundamentos del Código y regla de interpretación, de modo que todo se mira bajo otra perspectiva. \\[0.8em]

\textbf{¿Cómo se relacionan los derechos individuales y los de incidencia colectiva?} &
Existe una clara preeminencia del interés colectivo sobre el individual (arts. 14, 15 y 240 CCyCN). Por eso puede resultar necesario modelar el alcance del ejercicio de un derecho real en función de los intereses de la comunidad. \\

\midrule
\multicolumn{2}{@{}p{\textwidth}@{}}{%
\textbf{Resumen:} Los derechos reales son 14, creados exclusivamente por la ley federal y regidos por normas de orden público que los particulares no pueden modificar. Responden a un esquema de sociedad único para toda la Nación y se interpretan a la luz de la constitucionalización e internacionalización del derecho privado. El sistema consagra la propiedad privada, pero prevé la preeminencia del interés colectivo sobre el individual.
} \\
\end{longtable}

% ==========================================================
% CAPÍTULO 2
% ==========================================================
\chapter[Adquisición de los derechos reales]{Adquisición de los derechos reales}

\section{Formas de adquisición}

\subsection{Adquisición originaria y derivada}

Existen dos modalidades fundamentales de adquisición de un derecho real.

\begin{center}
\small
\renewcommand{\arraystretch}{1.3}
\begin{tabularx}{\textwidth}{@{}>{\raggedright\arraybackslash}p{0.17\textwidth}>{\raggedright\arraybackslash}X>{\raggedright\arraybackslash}p{0.30\textwidth}@{}}
\toprule
\textbf{Tipo} & \textbf{Concepto} & \textbf{Ejemplos} \\
\midrule
Originaria & El sistema reconoce la posibilidad de adquirir el derecho real en forma libre y absoluta, sin que derive de un vínculo jurídico anterior. & Cosas sin dueño (\textit{res nullius}); usucapión; adquisición legal. \\
Derivada & El derecho real es transmitido de un enajenante a un adquirente. Existe un vínculo jurídico previo. & Compraventa, donación, dación en pago, sucesión por causa de muerte. \\
\bottomrule
\end{tabularx}
\end{center}

\subsection{La adquisición derivada: actos entre vivos y causa de muerte}

La adquisición derivada puede producirse:

\begin{itemize}
    \item \textbf{Por actos entre vivos}: compraventa, permuta, donación, dación en pago, cesión, etc. Puede ser a título oneroso o gratuito.
    \item \textbf{Por causa de muerte}: el heredero ocupa el lugar del causante directamente y adquiere todo lo que aquel tenía.
\end{itemize}

\subsection{Regla \textit{nemo plus iuris} (art. 399)}

\begin{formula}{Nemo plus iuris (art. 399 CCyCN)}
Nadie puede transmitir a otro un derecho mejor o más extenso que el que tiene, y nadie puede adquirir sobre un objeto un derecho mejor o más extenso que el que tenía aquel de quien lo adquiere.
\end{formula}

Esta regla se aplica también a los derechos personales y, si bien admite excepciones, es importante tener presente que constituye la regla general.

\begin{formula}{Art. 1896 CCyCN}
El juez no puede constituir derechos reales ni imponer su constitución. (Principio de reserva legal en materia de derechos reales.)
\end{formula}

\section{Teoría del título y modo (art. 1892)}

La adquisición derivada de los derechos reales por actos entre vivos requiere necesariamente de dos elementos: el \textbf{título suficiente} y el \textbf{modo suficiente}.

\begin{formula}{Art. 1892 CCyCN}
El derecho real se adquiere con la concurrencia del título suficiente más el modo suficiente.
\end{formula}

\subsection{El título suficiente: concepto y alcance}

\begin{definition}{Título suficiente}
Es el acto jurídico revestido de las condiciones de fondo y de forma establecidas por la ley, que tiene por finalidad transmitir o constituir un derecho real.
\end{definition}

No debe identificarse el título en su sentido instrumental. Lo que se requiere para el derecho real es el título en su \textbf{sentido abstracto}, es decir, intangible. La causa del título puede ser una compraventa, una cesión, una permuta, una dación en pago o una donación, entre otras.

\begin{definition}{Acto jurídico (art. 259 CCyCN)}
Es el acto voluntario lícito que tiene por fin inmediato la adquisición, modificación o extinción de relaciones o situaciones jurídicas.
\end{definition}

El art. 260 CCyCN establece que el acto voluntario es el ejecutado con discernimiento, intención y libertad, que se manifiesta por un hecho exterior.

En materia de inmuebles, el título suficiente debe plasmarse en \textbf{escritura pública}.

% TODO: contenido incompleto o ambiguo en las notas originales (el relato de la analogía del docente es confuso en la transcripción)
\begin{example}{El título como acto intangible}
El docente ilustra el carácter abstracto del título con una analogía con el amor: no es posible decir <<te doy mi amor, dejámelo en el buzón>>. De igual modo, ante un problema de falta de título suficiente, no se resuelve afirmando que el título fue <<dejado en el buzón>>: lo que se había dejado era un instrumento que no plasmaba un acto jurídico causal. El título es intangible y abstracto; no puede <<dejarse en el buzón>>.
\end{example}

\subsection{El modo suficiente: la tradición posesoria}

El modo es el elemento material propio de los derechos reales que se ejercen por la posesión.

\begin{definition}{Tradición posesoria}
Es la entrega material de la cosa: una persona entrega materialmente la cosa a otra que la recibe. El poder de hecho que tenía el sujeto se transmite, con la entrega, a quien la adquiere.
\end{definition}

En la tradición, el sujeto que entrega es el \textbf{tradente} y el que recibe es el \textbf{accipiens}. La tradición es \textbf{constitutiva} en todos los casos en que se adquieren derechos reales que se ejercen por la posesión: sin tradición no se adquirió el derecho real.

\begin{formula}{Art. 1924 CCyCN}
La tradición no puede suplirse por fórmulas escritas. No puede transmitirse la posesión a través de cláusulas insertas en instrumentos escritos.
\end{formula}

\subsection{Requisitos de la tradición traslativa de dominio}

Para que la entrega transmita el dominio deben cumplirse los siguientes requisitos, que surgen del art. 1892 CCyCN.

\begin{center}
\small
\renewcommand{\arraystretch}{1.3}
\begin{tabularx}{\textwidth}{@{}>{\raggedright\arraybackslash}p{0.27\textwidth}>{\raggedright\arraybackslash}X@{}}
\toprule
\textbf{Requisito} & \textbf{Descripción} \\
\midrule
Legitimación & Debe ser realizada por el propietario. Solo el dueño puede transmitir el dominio (vinculado con el \textit{nemo plus iuris}). \\
Acto causal & Debe existir un título suficiente: compraventa, donación, dación en pago, etc. \\
Capacidad & Ambas partes deben ser capaces. \\
Escritura pública (inmuebles) & Si se trata de inmuebles, el acto jurídico causal debe plasmarse en escritura pública (art. 1017 CCyCN). \\
\bottomrule
\end{tabularx}
\end{center}

\subsection{La inscripción registral: efectos declarativos y constitutivos}

El efecto de la inscripción registral es distinto según el bien de que se trate.

\begin{center}
\small
\renewcommand{\arraystretch}{1.3}
\begin{tabularx}{\textwidth}{@{}>{\raggedright\arraybackslash}p{0.17\textwidth}>{\raggedright\arraybackslash}p{0.17\textwidth}>{\raggedright\arraybackslash}X@{}}
\toprule
\textbf{Tipo de bien} & \textbf{Efecto} & \textbf{Consecuencia práctica} \\
\midrule
Inmuebles & Declarativo & La inscripción hace oponible \textit{erga omnes} el derecho real. No es constitutiva del derecho, que ya nace con título más modo (tradición). \\
Automotores & Constitutivo & Sin transferencia (inscripción), no hay dominio. La inscripción es la que crea el derecho real. \\
\bottomrule
\end{tabularx}
\end{center}

\section{Excepciones a la tradición material}

La regla general es que la tradición debe ser material y efectiva, y no puede suplirse por ninguna fórmula escrita. Sin embargo, la ley admite dos excepciones.

\subsection{\textit{Traditio brevi manu} (arts. 1892 y 1933)}

Se produce cuando quien ya tiene la cosa como \textbf{tenedor} (por ejemplo, en calidad de locatario) pasa a ser \textbf{titular del derecho real} (propietario) por un acto jurídico causal. La relación de poder \textbf{asciende de categoría}: el sujeto pasa de una relación de poder inferior (tenencia) a una superior (dominio).

\begin{example}{María y el departamento (\textit{traditio brevi manu})}
María, mujer joven y docente, alquiló durante muchos años un departamento. En un momento se le presentó la posibilidad de comprarlo al propietario, y realizó el acto jurídico causal (la compraventa, plasmada en escritura pública). Sin la \textit{traditio brevi manu}, María debería contratar una mudadora, cargar sus muebles, salir del departamento en calidad de inquilina y esperar que ingrese el propietario para que luego se lo entregue como propietaria. Semejante dispendio material es innecesario: la ley admite que, haciéndolo constar en la escritura, María ascienda de categoría (de tenedora o locataria a propietaria) sin necesidad de entrega material.
\end{example}

\subsection{Constituto posesorio (arts. 1892 y 1923)}

Es el caso inverso: quien tenía la cosa como \textbf{propietario} la transmite a un tercero y pasa a tenerla como \textbf{tenedor} (locatario, comodatario, etc.). La relación de poder \textbf{desciende de categoría}.

% TODO: contenido incompleto o ambiguo en las notas originales (el diálogo del ejemplo está truncado en la transcripción)
\begin{example}{El propietario endeudado (constituto posesorio)}
Un señor que en el año 2001 tenía muchas deudas decide vender su departamento para pagarlas. Lo que cambia es la calidad en que se encuentra en relación con la cosa: antes era propietario y ahora será locatario. Sin el constituto posesorio, debería entregar materialmente la posesión al adquirente y luego volver a ingresar como inquilino. El constituto posesorio evita ese dispendio: el propietario que vende permanece en el inmueble, pero ahora en calidad de locatario del nuevo dueño, haciéndolo constar con la debida constancia notarial y fecha cierta.
\end{example}

\subsection{Cuadro comparativo de las excepciones}

\begin{center}
\small
\renewcommand{\arraystretch}{1.3}
\begin{tabularx}{\textwidth}{@{}>{\raggedright\arraybackslash}p{0.22\textwidth}>{\raggedright\arraybackslash}X>{\raggedright\arraybackslash}X>{\raggedright\arraybackslash}p{0.16\textwidth}@{}}
\toprule
\textbf{Instituto} & \textbf{Situación previa} & \textbf{Situación posterior} & \textbf{Movimiento} \\
\midrule
\textit{Traditio brevi manu} (arts. 1892 y 1933) & Tenedor (locatario, comodatario) & Propietario (titular del dominio) & Ascendente \\
Constituto posesorio (arts. 1892 y 1923) & Propietario (titular del dominio) & Tenedor (locatario, comodatario) & Descendente \\
\bottomrule
\end{tabularx}
\end{center}

\section{Adquisición legal u originaria por disposición de la ley (art. 1894)}

Existen casos en que la ley misma decide que el derecho se adquiere, sin necesidad de título ni modo. El art. 1894 CCyCN establece los casos de adquisición legal u originaria: el derecho real de habitación del cónyuge supérstite, el del conviviente supérstite y los condominios con indivisión forzosa.

\subsection{Derecho real de habitación del cónyuge supérstite}

Al fallecer uno de los cónyuges, la ley reconoce al sobreviviente el derecho real de habitación sobre el inmueble que era asiento del hogar conyugal, con carácter gratuito y vitalicio, siempre que se cumplan los requisitos legales. La ley misma constituye el derecho.

\subsection{Derecho real del conviviente supérstite}

El nuevo Código incorporó al conviviente (unión convivencial) entre los beneficiarios de este derecho. Esto se presenta como una manifestación del cambio de paradigma del nuevo cuerpo normativo: <<el conviviente impensado>>.

\subsection{Condominios con indivisión forzosa}

Es otra forma de adquisición legal u originaria, que será analizada en profundidad al tratar cada derecho real en particular.

\section{Cuadro Cornell del capítulo}

\begin{longtable}{@{}>{\raggedright\arraybackslash}p{0.27\textwidth}>{\raggedright\arraybackslash}p{0.67\textwidth}@{}}
\toprule
\textbf{Preguntas / palabras clave} & \textbf{Notas / respuestas} \\
\midrule
\endfirsthead
\toprule
\textbf{Preguntas / palabras clave} & \textbf{Notas / respuestas} \\
\midrule
\endhead
\bottomrule
\endfoot

\textbf{¿Cuáles son las formas de adquisición de los derechos reales?} &
Originaria: se adquiere sin vínculo jurídico anterior (cosas sin dueño, usucapión, adquisición legal). Derivada: se transmite de un enajenante a un adquirente. La derivada puede ser por actos entre vivos (onerosos o gratuitos) o por causa de muerte. \\[0.8em]

\textbf{¿Qué es el \textit{nemo plus iuris}?} &
Es la regla del art. 399 CCyCN: nadie puede transmitir a otro un derecho mejor o más extenso que el que tiene. Se aplica también a los derechos personales y, aunque admite excepciones, es la regla general. \\[0.8em]

\textbf{¿En qué consiste la teoría del título y modo?} &
Para adquirir un derecho real de modo derivado por actos entre vivos se requieren dos elementos concurrentes (art. 1892 CCyCN). 1) Título suficiente: el acto jurídico causal (compraventa, donación, etc.), entendido en sentido abstracto e intangible; en inmuebles debe constar en escritura pública. 2) Modo suficiente: la tradición posesoria, es decir, la entrega material de la cosa, que es constitutiva y no puede suplirse por fórmulas escritas (art. 1924 CCyCN). \\[0.8em]

\textbf{¿Qué requisitos tiene la tradición traslativa de dominio?} &
Legitimación (la realiza el propietario), acto causal (título suficiente), capacidad de ambas partes y, en inmuebles, escritura pública (art. 1017 CCyCN). \\[0.8em]

\textbf{¿Cuál es el efecto de la inscripción registral?} &
En inmuebles es declarativo: hace oponible el derecho \textit{erga omnes}, pero no lo constituye, pues el derecho ya nació con título más modo. En automotores es constitutivo: sin inscripción no hay dominio. \\[0.8em]

\textbf{¿Qué son la \textit{traditio brevi manu} y el constituto posesorio?} &
Son las dos excepciones a la tradición material. En la \textit{traditio brevi manu} (arts. 1892 y 1933) el tenedor pasa a ser propietario sin entrega material (ascenso de categoría). En el constituto posesorio (arts. 1892 y 1923) el propietario pasa a ser tenedor sin devolución física de la cosa (descenso de categoría). En el constituto posesorio se hace constar con constancia notarial y fecha cierta. \\[0.8em]

\textbf{¿Cuáles son los casos de adquisición legal de derechos reales?} &
El art. 1894 CCyCN establece los casos en que la ley misma constituye el derecho real: habitación del cónyuge supérstite, habitación del conviviente supérstite (novedad del nuevo Código) y condominios con indivisión forzosa. \\

\midrule
\multicolumn{2}{@{}p{\textwidth}@{}}{%
\textbf{Resumen:} La adquisición puede ser originaria o derivada, y la derivada puede darse por actos entre vivos o por causa de muerte, con el límite del \textit{nemo plus iuris}. Por actos entre vivos se exige título suficiente (acto jurídico causal, con escritura pública en inmuebles) y modo suficiente (tradición material, constitutiva), salvo las excepciones de la \textit{traditio brevi manu} y el constituto posesorio. La inscripción es declarativa en inmuebles y constitutiva en automotores. La ley también constituye directamente ciertos derechos reales (art. 1894 CCyCN).
} \\
\end{longtable}

% ==========================================================
% CAPÍTULO 3
% ==========================================================
\chapter[Publicidad, transmisión y extinción]{Publicidad, transmisión y extinción de los derechos reales}

\section{Publicidad de los derechos reales (art. 1893)}

Los derechos reales deben ser respetados por todos y, si deben ser respetados por todos, \textbf{deben ser conocidos por todos}. De allí el principio de publicidad.

\begin{formula}{Art. 1893 CCyCN}
Norma que regula la publicidad de los derechos reales y la oponibilidad a terceros de las relaciones de poder.
\end{formula}

\subsection{Publicidad posesoria}

La forma más antigua y natural de publicidad es la posesión misma: desde la existencia del hombre en sociedad, el acto más claro de un vínculo con la cosa es tenerla en posesión. Esa es la \textbf{publicidad posesoria}.

\subsection{Publicidad registral}

La inscripción en un registro público permite que los terceros conozcan quién es el titular del derecho real sobre determinados bienes. El efecto de la inscripción puede ser declarativo (inmuebles) o constitutivo (automotores), según se desarrolló en el capítulo anterior.

\section{Transmisión de los derechos reales}

Todos los derechos, tanto reales como personales, son en principio transmisibles. La transmisión opera bajo las mismas reglas ya vistas: si el derecho se adquiere con título suficiente y modo, quien lo transmite lo hizo con título suficiente y modo, y quien lo transmite a su vez debe hacerlo de esa misma manera.

\begin{important}
La regla \textit{nemo plus iuris} (art. 399 CCyCN) es el límite infranqueable de la transmisión: nadie puede transmitir más derechos de los que tiene.
\end{important}

\section{Extinción de los derechos reales (art. 1907)}

\begin{formula}{Art. 1907 CCyCN}
El derecho real se extingue en forma relativa o absoluta.
\end{formula}

\subsection{Extinción relativa}

En la extinción relativa el derecho real se extingue para el titular actual pero continúa en otro sujeto. Ocurre cuando el titular dona, vende o entrega la cosa en pago: pierde el dominio, pero la cosa y el derecho sobre ella continúan, ahora en cabeza del adquirente.

\subsection{Extinción absoluta}

En la extinción absoluta el derecho real desaparece por completo, generalmente por destrucción de la cosa. Como el derecho real es directamente un vínculo entre el titular y la cosa, cuando la cosa se destruye por cualquier medio ese vínculo se extingue necesariamente.

% TODO: contenido incompleto o ambiguo en las notas originales (la transcripción dice <<destrucción de la causa>>; se interpreta <<cosa>> por el contexto de la propia explicación)

\subsection{La consolidación}

Es una forma particular de extinción, que ocurre cuando se reúnen en una misma persona el derecho real principal y el desmembrado: por ejemplo, el nudo propietario adquiere el usufructo, o el usufructuario adquiere el dominio pleno. Es análoga a lo que en los derechos personales se denomina \textbf{confusión}.

\subsection{Cuadro comparativo de las formas de extinción}

\begin{center}
\small
\renewcommand{\arraystretch}{1.3}
\begin{tabularx}{\textwidth}{@{}>{\raggedright\arraybackslash}p{0.15\textwidth}>{\raggedright\arraybackslash}X>{\raggedright\arraybackslash}p{0.17\textwidth}>{\raggedright\arraybackslash}p{0.22\textwidth}@{}}
\toprule
\textbf{Tipo} & \textbf{Descripción} & \textbf{¿El derecho continúa?} & \textbf{Ejemplo} \\
\midrule
Relativa & Se extingue para el titular actual. & Sí, en otro sujeto. & Venta, donación, dación en pago. \\
Absoluta & El derecho desaparece por completo. & No. & Destrucción de la cosa. \\
Consolidación & Se reúnen en un mismo sujeto el derecho real y el que lo desmembraba. & Se consolida en uno solo. & El usufructuario adquiere el dominio pleno. \\
\bottomrule
\end{tabularx}
\end{center}

\section{Reflexión final: el derecho real como relación de poder}

El derecho real es una \textbf{relación de poder del titular respecto de la cosa}. Ese esquema (esa línea o flecha) es lo que permite perseguir la cosa, en ocasiones, como ocurre en las hipotecas, aun cuando hubiera cambiado de titular. Esa relación de poder es defendida y protegida por el sistema jurídico y debe ser respetada por toda la sociedad. De allí la indefectible necesidad de la publicidad, ya sea posesoria o registral, a pesar de los límites que la Nación decida imponer en función de los intereses de la comunidad.

Desde el nacimiento y a lo largo de la existencia, todos los seres humanos tienen contactos permanentes con las cosas. Ese contacto directo e inmediato tiene una protección jurídica que debe ser respetada, porque es una manera de respetar también la paz social.

En el próximo tema se tratarán todos los vínculos de poder, y se verá que no siempre están en cabeza del titular del derecho real: no siempre las personas viven en las tierras de las que son propietarias, y existen diversas relaciones de poder.

\section{Cuadro Cornell del capítulo}

\begin{longtable}{@{}>{\raggedright\arraybackslash}p{0.27\textwidth}>{\raggedright\arraybackslash}p{0.67\textwidth}@{}}
\toprule
\textbf{Preguntas / palabras clave} & \textbf{Notas / respuestas} \\
\midrule
\endfirsthead
\toprule
\textbf{Preguntas / palabras clave} & \textbf{Notas / respuestas} \\
\midrule
\endhead
\bottomrule
\endfoot

\textbf{¿Por qué los derechos reales requieren publicidad?} &
Porque son oponibles \textit{erga omnes}: todos deben respetarlos y, para ello, deben conocerlos. La publicidad puede ser posesoria (la posesión misma como hecho visible) o registral (inscripción en registros públicos, con efectos declarativos o constitutivos según el bien). \\[0.8em]

\textbf{¿Cómo opera la transmisión de los derechos reales?} &
Todos los derechos son en principio transmisibles. La transmisión sigue las mismas reglas de la adquisición (título suficiente y modo), con el \textit{nemo plus iuris} (art. 399 CCyCN) como límite: nadie puede transmitir más derechos de los que tiene. \\[0.8em]

\textbf{¿Cómo se extingue el derecho real?} &
Según el art. 1907 CCyCN, la extinción puede ser relativa (el derecho se extingue para el titular actual pero continúa en otro sujeto: venta, donación) o absoluta (el derecho desaparece por completo, generalmente por destrucción de la cosa). \\[0.8em]

\textbf{¿Qué es la consolidación?} &
Es la reunión en una misma persona del derecho real principal y del que lo desmembraba (por ejemplo, el usufructuario adquiere el dominio pleno). Es análoga a la confusión en los derechos personales. \\[0.8em]

\textbf{¿Qué es el derecho real en esencia?} &
Es una relación de poder del titular respecto de la cosa, defendida y protegida por el sistema jurídico y que debe ser respetada por toda la sociedad. Permite perseguir la cosa, como en las hipotecas, aun cuando hubiera cambiado de titular, y por eso requiere publicidad posesoria o registral. \\

\midrule
\multicolumn{2}{@{}p{\textwidth}@{}}{%
\textbf{Resumen:} Los derechos reales deben ser conocidos por todos mediante la publicidad posesoria o registral, porque son relaciones de poder oponibles a la sociedad. Se transmiten con título y modo, con el límite del \textit{nemo plus iuris}, y se extinguen de forma relativa o absoluta; la consolidación es una forma particular de extinción. Todo ello se enmarca en la constitucionalización del derecho privado y en la tensión entre los derechos individuales y los de incidencia colectiva, siendo el derecho real, en esencia, una relación de poder del titular sobre la cosa.
} \\
\end{longtable}

\end{document}
